\chapter{Introduction}
\label{chap:intro}

\section{General Presentation of the Internship}
This report was prepared as part of a Master's internship carried out at \textbf{Outsight}, in \textbf{Sophia Antipolis, Valbonne}, France, between \textbf{16 February 2026 and 31 July 2026}. The internship was completed within Outsight's engineering team, and consisted in analyzing and carrying out the migration of Outsight Cloud (OC), the company's platform for managing LiDAR-related workflows, simulations, organizations, sites, and associated resources, from a legacy virtual-machine-based infrastructure toward a cloud-native environment on AWS. Chapter~\ref{chap:company} presents the company and its business sector, Chapter~\ref{chap:internship-topic} presents the internship topic in more detail, and Section~\ref{sec:report-structure} outlines the remainder of this report.

\section{Industrial Context and Legacy Constraints}
Cloud computing has transformed the way software systems are designed, deployed, and maintained, offering scalable infrastructure, managed services, and automated deployment capabilities that are difficult to achieve with self-managed environments. As a consequence, many organizations are progressively modernizing legacy applications toward cloud-native architectures.

Migrating an existing production platform to the cloud, however, is rarely straightforward. Legacy systems are typically the result of years of incremental development shaped by evolving business requirements, and over time they accumulate architectural dependencies and operational practices that complicate scalability, maintainability, and long-term evolution.

This report analyzes the migration of Outsight Cloud (OC), a platform used to manage LiDAR-related workflows, simulations, organizations, sites, and associated resources. Before the migration, the platform relied on a virtual-machine-based infrastructure hosted on DigitalOcean Droplets and orchestrated through Docker Compose, integrating several external services, most notably Airtable for core data management. This architecture enabled rapid development during the early product lifecycle, but as usage and integrations expanded it began to show clear signs of architectural strain.

\section{Report Structure}
\label{sec:report-structure}
This report is organized as follows:

\begin{itemize}
    \item \textbf{Chapter~\ref{chap:company}} presents the company, Outsight, and its business sector.
    \item \textbf{Chapter~\ref{chap:internship-topic}} presents the internship topic, i.e.\ the engineering problem addressed in this work.
    \item \textbf{Chapter~\ref{chap:background}} introduces the conceptual and related-work foundations on cloud-native modernization, migration patterns, and IAM evolution in legacy systems.
    \item \textbf{Chapter~\ref{chap:baseline}} describes the case-study baseline (as-is architecture), including technical debt and cost-related migration drivers.
    \item \textbf{Chapter~\ref{chap:methodology}} presents the migration methodology and decision framework, including candidate options and selection rationale.
    \item \textbf{Chapter~\ref{chap:architecture}} defines the target architecture (to-be), its main components, and design trade-offs.
    \item \textbf{Chapter~\ref{chap:implementation}} details the incremental migration implementation, with focus on IaC, data migration tooling, and the validation pipeline.
    \item \textbf{Chapter~\ref{chap:evaluation}} discusses the evaluation approach and the evidence on migration cost, performance, and operational impact.
    \item \textbf{Chapter~\ref{chap:conclusions}} concludes by summarizing contributions and outlining future work.
\end{itemize}
